\subsection{EXISTENCE OF PRIMITIVES}

Let f be a function defined on an arbitrary interval \(I \subset R\) ; for a function g defined on I to be a primitive of f, it is necessary and sufficient that the restriction of g to every compact interval \(J \subset I\) be a primitive of the restriction of f to J.

THEOREM 1. Let A be a set filtered by a filter \(\mathfrak{F}\) , and \((\mathbf{f}_{\alpha})_{\alpha \in \mathrm{A}}\) a family of vector functions with values in a complete normed space E over \(\mathbf{R}\) , defined on an interval \(\mathbf{I} \subset \mathbf{R}\) : for each \(\alpha \in \mathrm{A}\) let \(\mathbf{g}_{\alpha}\) be a primitive of \(\mathbf{f}_{\alpha}\) . We suppose that:

\(1^{\circ}\) with respect to the filter \(\mathfrak{F}\) the functions \(\mathbf{f}_{\alpha}\) converge uniformly on every compact subset of I to a function \(\mathbf{f}\) ;

\(2^{\circ}\) there is a point \(a \in I\) such that, with respect to the filter \(\mathfrak{F}\) , the family \((\mathbf{g}_{\alpha}(a))\) has a limit in E.

Under these hypotheses the functions \(g_{\alpha}\) converge uniformly (with respect to F) on every compact subset of I to a primitive g of f.

By the remark at the beginning of this subsection we can restrict ourselves to the case where I is a compact interval.

First let us show that the \(g_{\alpha}\) converge uniformly on I to a continuous function g. By hypothesis, for every \(\varepsilon > 0\) there is a set \(M \in F\) such that, for any two indices \(\alpha, \beta\) belonging to M, one has \(\left\|\mathbf{f}_{\alpha}(x) - \mathbf{f}_{\beta}(x)\right\| \leqslant \varepsilon\) for each \(x \in I\) ; in consequence one has (I, p.~15, th. 2)

\[
\left\| \mathbf {g} _ {\alpha} (x) - \mathbf {g} _ {\beta} (x) - (\mathbf {g} _ {\alpha} (a) - \mathbf {g} _ {\beta} (a)) \right\| \leqslant \varepsilon | x - a | \leqslant \varepsilon l
\]

where \(l\) denotes the length of I; since by hypothesis \(\mathbf{g}_{\alpha}(a)\) approaches a limit with respect to \(\mathfrak{F}\) , it follows from the Cauchy criterion that the \(\mathbf{g}_{\alpha}\) converge uniformly on I. It remains to show that the limit \(\mathbf{g}\) of the \(\mathbf{g}_{\alpha}\) is a primitive of \(\mathbf{f}\) .

For each integer n \textgreater{} 0 let \(\alpha_{n}\) be an index such that \(\left\|\mathbf{f}(x) - \mathbf{f}_{\alpha_{n}}(x)\right\| \leqslant 1/n\) on I; it is clear that the sequence \(\left(\mathbf{f}_{\alpha_{n}}\right)\) converges uniformly to f and that the sequence \(\left(\mathbf{g}_{\alpha_{n}}\right)\) converges uniformly to g on I. Let \(H_{n}\) be the countable subset of I where \(f_{\alpha_{n}}\) is not the derivative of \(g_{\alpha_{n}}\) , and let H be the union of the \(H_{n}\) , which is thus a countable subset of I; we shall see that at every point \(x \in \mathrm{I}\) not belonging to \(\mathrm{H}\) the function \(\mathbf{g}\) has a derivative equal to \(\mathbf{f}(x)\) . Indeed, one sees as above that for every \(m \geqslant n\) and every \(y \in \mathrm{I}\) one has

\[
\left| \left| \mathbf {g} _ {\alpha_ {m}} (y) - \mathbf {g} _ {\alpha_ {m}} (x) - \left(\mathbf {g} _ {\alpha_ {n}} (y) - \mathbf {g} _ {\alpha_ {n}} (x)\right) \right| \right| \leqslant \frac {2}{n} | y - x |.
\]

Letting \(m\) increase indefinitely one also has

\[
\left| \left| \mathbf {g} (y) - \mathbf {g} (x) - \left(\mathbf {g} _ {\alpha_ {n}} (y) - \mathbf {g} _ {\alpha_ {n}} (x)\right) \right| \right| \leqslant \frac {2}{n} | y - x |
\]

for every \(y \in \mathbf{I}\) ; now, there exists an \(h > 0\) such that, for \(|y - x| \leqslant h\) and \(y \in \mathbf{I}\) , one has \(\left\| \mathbf{g}_{\alpha}(y) - \mathbf{g}_{\alpha_n}(x) - \mathbf{f}_{\alpha_n}(x)(y - x) \right\| \leqslant |y - x| / n\) ; since, on the other hand, we have \(\left\| \mathbf{f}(x) - \mathbf{f}_{\alpha_n}(x) \right\| \leqslant 1 / n\) , we finally obtain

\[
\| \mathbf {g} (y) - \mathbf {g} (x) - \mathbf {f} (x) (y - x) \| \leqslant \frac {4}{n} | y - x |
\]

for \(y \in I\) and \(|y - x| \leqslant h\) , which completes the proof.

COROLLARY 1. The set H of maps from I into E which admit a primitive on an interval I is a closed (and so complete) vector subspace of the complete vector space \(\mathcal{F}_{c}(\mathrm{I};\mathrm{E})\) of maps from I into E, endowed with the topology of uniform convergence on every compact subset of I (Gen.~Top., X, p.~277).

COROLLARY 2. Let \(x_0\) be a point of I, and for each function \(\mathbf{f} \in \mathcal{H}\) let \(\mathrm{P}(\mathbf{f})\) be the primitive of \(\mathbf{f}\) which vanishes at the point \(x_0\) ; the map \(\mathbf{f} \mapsto \mathrm{P}(\mathbf{f})\) of \(\mathcal{H}\) into \(\mathcal{F}_c(\mathrm{I};\mathrm{E})\) is a continuous linear mapping.

Cor. 1 to th. 1 allows us to establish the existence of primitives for certain categories of functions by the following procedure: if one knows that the functions belonging to a subset A of \(\mathcal{F}_{c}(\mathrm{I};\mathrm{E})\) admit a primitive, so will the functions belonging to the closure in \(\mathcal{F}_{c}(\mathrm{I};\mathrm{E})\) of the vector subspace generated by A. We shall apply this method in the next subsection.

